\frontmatter

% ============================================================
% Halaman Sampul
% ============================================================
\begin{titlepage}
  \centering
  \vspace*{1cm}
  {\LARGE\bfseries Pemrograman Web\par}
  \vspace{0.5cm}
  {\large Buku Ajar Berbasis Outcome-Based Education (OBE)\par}
  \vspace{1cm}
  {\large Mata Kuliah: Pemrograman Web\par}
  {\large Kode: TIF-XXX\par}
  \vspace{2cm}
  {\large Program Studi Teknik Informatika\par}
  {\large Fakultas Teknologi Informasi\par}
  {\large Universitas Bale Bandung\par}
  \vfill
  {\large 2026\par}
\end{titlepage}

\cleardoublepage

% ============================================================
% Kata Pengantar
% ============================================================
\chapter*{Kata Pengantar}
\addcontentsline{toc}{chapter}{Kata Pengantar}

Puji syukur ke hadirat Tuhan Yang Maha Esa atas tersusunnya buku ajar \textit{Pemrograman Web} ini. Buku ini disusun dengan pendekatan \textbf{Outcome-Based Education (OBE)} agar capaian pembelajaran terukur dan selaras dengan Capaian Pembelajaran Lulusan (CPL) serta Capaian Pembelajaran Mata Kuliah (CPMK).

Pemrograman web merupakan keterampilan inti bagi mahasiswa Teknik Informatika. Materi di dalam buku ini memandu Anda dari pengenalan kebutuhan pengguna, pembangunan antarmuka, hingga pengembangan back-end, integrasi data, keamanan, serta deployment.

Buku ini dirancang untuk pembelajar pemula, sehingga setiap bagian dilengkapi contoh, latihan, dan studi kasus sederhana. Semoga buku ini membantu Anda membangun fondasi yang kuat dan menerapkan praktik terbaik dalam pengembangan web.

\vspace{1cm}
\begin{flushright}
Penyusun\\
TBD
\end{flushright}

\cleardoublepage

% ============================================================
% Cara Menggunakan Buku Ini
% ============================================================
\chapter*{Cara Menggunakan Buku Ini}
\addcontentsline{toc}{chapter}{Cara Menggunakan Buku Ini}

Buku ajar ini dirancang dengan pendekatan OBE untuk memaksimalkan pencapaian pembelajaran Anda. Berikut panduan penggunaan buku ini:

\section*{Struktur Buku}

\textbf{Bab I: Pendahuluan dan Orientasi}\\
Memperkenalkan tujuan buku, keterkaitan dengan RPS, dan konteks kurikulum OBE.

\textbf{Bab II: Landasan Teori}\\
Menyajikan fondasi konseptual pemrograman web dan arsitektur client-server. Bab ini mendukung \textbf{Pertemuan 1--2} RPS (pengantar, kebutuhan pengguna, arsitektur) serta menjadi prasyarat konseptual untuk \textbf{Pertemuan 4--7} (HTML, CSS).

\textbf{Bab III–XV: Unit Materi Inti}\\
Setiap bab mencakup satu topik utama pemrograman web dengan struktur lengkap: Sub-CPMK, materi, aktivitas, latihan, asesmen, dan checklist.

\textbf{Bab XVI: Evaluasi, Refleksi, dan Integrasi Kompetensi}\\
Berisi kumpulan asesmen (Quiz 1--4, Tugas, UTS, UAS, Project), rubrik penilaian komprehensif, tinjauan pencapaian kompetensi, dan rekomendasi studi lanjutan.

\textbf{Lampiran}\\
Menyediakan rubrik penilaian, template tugas, glosarium, dan referensi tambahan.

\section*{Komponen dalam Setiap Bab}

\begin{enumerate}
  \item \textbf{Sub-CPMK}: Kompetensi spesifik yang harus dicapai pada bab terkait.
  \item \textbf{Materi Pokok}: Penjelasan konsep disertai contoh dan ilustrasi.
  \item \textbf{Aktivitas Pembelajaran}: Tugas, praktik, atau studi kasus.
  \item \textbf{Latihan}: Soal atau refleksi untuk menguatkan pemahaman.
  \item \textbf{Asesmen}: Instrumen untuk mengukur capaian.
  \item \textbf{Checklist}: Daftar indikator yang harus dikuasai.
\end{enumerate}

\section*{Tips Belajar Efektif}

\begin{itemize}
  \item Praktikkan setiap contoh kode dan modifikasi agar memahami dampaknya.
  \item Catat istilah baru dan cek kembali di glosarium.
  \item Kerjakan latihan sebelum melihat solusi.
  \item Gunakan project mini sebagai latihan integrasi.
  \item Diskusikan kesulitan dengan dosen atau teman.
\end{itemize}

\cleardoublepage

% ============================================================
% Identitas Mata Kuliah
% ============================================================
\chapter*{Identitas Mata Kuliah}
\addcontentsline{toc}{chapter}{Identitas Mata Kuliah}

\begin{tabular}{ll}
  Nama Program Studi & : Teknik Informatika \\
  Nama Mata Kuliah & : Pemrograman Web \\
  Kode Mata Kuliah & : TIF-XXX \\
  Semester & : 3 (Tiga) \\
  SKS / Bobot Kredit & : 2 SKS (70\% Praktik, 30\% Teori) \\
  Dosen Pengampu & : TBD \\
  Tanggal Penyusunan & : 2 Maret 2026 \\
\end{tabular}

\vspace{1cm}

\section*{Capaian Pembelajaran Lulusan (CPL)}

CPL yang dibebankan pada mata kuliah ini mencakup kompetensi lulusan dalam aspek pengetahuan, keterampilan, dan sikap:

\begin{enumerate}
  \item \textbf{CPL-04:} Kemampuan mendesain, mengimplementasi dan mengevaluasi solusi berbasis computing yang memenuhi kebutuhan-kebutuhan computing pada sebuah disiplin program.
  \item \textbf{CPL-05:} Mengidentifikasi dan menganalisis kebutuhan-kebutuhan pengguna dan mempertimbangkannya dalam memilih, membuat, mengintegrasi, mengevaluasi, dan mengadministrasi sistem berbasis computing.
\end{enumerate}

\section*{Capaian Pembelajaran Mata Kuliah (CPMK)}

Kemampuan atau kompetensi spesifik yang diharapkan mahasiswa kuasai setelah menyelesaikan mata kuliah:

\begin{enumerate}
  \item \textbf{CPMK-1:} Analisis kebutuhan pengguna dan perancangan antarmuka (UI/UX).
  \item \textbf{CPMK-2:} Implementasi solusi web front-end (HTML, CSS, JavaScript).
  \item \textbf{CPMK-3:} Pengembangan sistem back-end dan integrasi data (PHP, MySQL, API).
  \item \textbf{CPMK-4:} Evaluasi, deployment, dan administrasi web (testing, security, release).
\end{enumerate}

\section*{Matriks Keterkaitan CPL dan CPMK}

\begin{table}[htbp]
\centering
\begin{tabular}{|c|c|c|P{6cm}|}
  \hline
  \textbf{CPL} & \textbf{CPMK} & \textbf{Kontribusi} & \textbf{Keterangan} \\
  \hline
  CPL-04 & CPMK-2 & Tinggi & Implementasi antarmuka web responsif dan dinamis \\
  \hline
  CPL-04 & CPMK-3 & Tinggi & Pengembangan back-end dan integrasi data \\
  \hline
  CPL-04 & CPMK-4 & Sedang & Evaluasi performa dan deployment \\
  \hline
  CPL-05 & CPMK-1 & Tinggi & Analisis kebutuhan pengguna dan UI/UX \\
  \hline
  CPL-05 & CPMK-3 & Sedang & Integrasi kebutuhan pengguna ke sistem \\
  \hline
  CPL-05 & CPMK-4 & Sedang & Administrasi dan evaluasi sistem \\
  \hline
\end{tabular}
\caption{Matriks Keterkaitan CPL dan CPMK}
\end{table}

\cleardoublepage

% ============================================================
% Daftar Isi dan Daftar Pendukung
% ============================================================
\phantomsection
\addcontentsline{toc}{chapter}{Daftar Isi}
\tableofcontents
\cleardoublepage

\clearpage
\phantomsection
\addcontentsline{toc}{chapter}{Daftar Gambar}
\listoffigures
\cleardoublepage

\phantomsection
\addcontentsline{toc}{chapter}{Daftar Tabel}
\listoftables
\cleardoublepage

\phantomsection
\addcontentsline{toc}{chapter}{Daftar Kode Program}
\lstlistoflistings
\cleardoublepage

